\documentclass[journal]{IEEEtran}

% =========================================================================
% PACKAGES & CONFIGURATION
% =========================================================================
\usepackage{amsmath,amssymb,amsfonts,bm}
\usepackage{graphicx}
\usepackage{booktabs}
\usepackage{multirow}
\usepackage{array}
\usepackage{makecell}
\usepackage{cite}
\usepackage{url}
\usepackage{microtype}
\usepackage{balance}
\usepackage{xcolor}
\usepackage[pagebackref=false,breaklinks=true,colorlinks=true,citecolor=blue,linkcolor=blue,urlcolor=blue]{hyperref}

% Correct bad hyphenation here
\hyphenation{op-ti-cal net-works semi-conduc-tor dis-en-tan-gle-ment pa-thol-o-gy}

\begin{document}

% =========================================================================
% TITLE & AUTHOR BLOCK
% =========================================================================
\title{DFU-WAGNER net: Multi-Task Diabetic Foot Ulcer Characterization via Hybrid Feature Fusion, Subspace Disentanglement, and Uncertainty-Gated Domain Alignment}

\author{
    Priyanka~Saxena,~
    Aniketh~Vuruputuri,~
    Nikunj~Jayesh,~
    U.~V.~S.~Pavan~Kumar,~
    Aditi,~
    and~Jagadhabhiram~Medishetty%
    \thanks{P. Saxena, A. Vuruputuri, N. Jayesh, U. V. S. Pavan Kumar, Aditi, and J. Medishetty are with the Department of Computer Science and Engineering / Department of Biomedical Engineering (e-mail: corresponding.author@research.org).}%
}

% Paper headers
\markboth{IEEE Transactions on Medical Imaging,~Vol.~XX, No.~X, October~2026}%
{Saxena \MakeLowercase{\textit{et al.}}: DFU-WAGNER net: Multi-Task Diabetic Foot Ulcer Characterization}

\maketitle

% =========================================================================
% ABSTRACT
% =========================================================================
\begin{abstract}
Diabetic foot ulcers (DFUs) represent a severe global healthcare burden requiring concurrent assessment along two clinically orthogonal diagnostic axes: microbiological/etiological pathology (\textit{None}, \textit{Infection}, \textit{Ischaemia}, \textit{Both}) and anatomical tissue invasion depth (Wagner Grades 0--4). However, automated diagnostic algorithms face critical clinical and computational bottlenecks: existing public cohorts are strictly fragmented---annotating either pathology or depth, but never both; severe class imbalance suppresses limb-threatening ischaemia detection (3.8\% prevalence); and cross-cohort acquisition differences induce negative transfer across disjoint hospital sources. 

To overcome these challenges, this paper presents \textbf{DFU-WAGNER net}, a unified multi-task deep representation learning framework. DFU-WAGNER net couples an \textbf{EfficientNetV2-S} convolutional stream with a \textbf{Swin-Transformer Tiny} stream through symmetrical 1,280-dimensional linear projections, capturing localized epidermal textures alongside macroscopic spatial foot topology within a fused 2,560-dimensional representation. Rather than utilizing monolithic feature sharing, a dedicated three-way \textbf{TaskDisentangler} partitions the latent space into an $L_2$-normalized shared anatomical subspace ($z_{\text{shared}} \in \mathbb{R}^{768}$) and two task-isolated private subspaces ($z_{\text{path\_private}} \in \mathbb{R}^{512}$ and $z_{\text{depth\_private}} \in \mathbb{R}^{512}$), strictly constrained by a pairwise cosine \textbf{Orthogonality Regularization Loss} ($\mathcal{L}_{\text{orth}}$). 

To prevent cross-cohort distribution collapse, we introduce \textbf{Uncertainty-Gated Maximum Mean Discrepancy (UG-MMD)}, which aligns \textit{only} the shared subspace between the DFUC2021 cohort and the Wagner cohort using multi-scale Gaussian kernels dynamically modulated by detached predictive entropy confidence weights. Long-tailed pathology classification is optimized via \textbf{Class-Balanced Focal Loss} and \textbf{Hard-Pair Tangent LMCoT Metric Learning} on $z_{\text{shared}}$, while anatomical ulcer depth is modeled using \textbf{Consistent Rank Logits (CORAL)} ordinal cumulative cutoffs. 

The primary 59,783,374-parameter architecture was evaluated across the DFUC2021 benchmark ($n = 5,955$ labeled samples) and a non-destructively deduplicated Wagner cohort ($n = 11,491$ training samples). Comprehensive experimental evaluation demonstrates that DFU-WAGNER net achieves a state-of-the-art Macro-F1 of \textbf{0.9201} and an overall accuracy of \textbf{90.18\%} on DFUC2021, while decisively resolving the extreme minority bottleneck by elevating \textit{Ischaemia} F1-score to \textbf{0.9149} with \textbf{93.48\%} sensitivity. Concurrently, the network achieves an ordinal Quadratic Weighted Kappa (QWK) of \textbf{0.7765}, adjacent depth accuracy of \textbf{91.42\%}, and Mean Absolute Error (MAE) of \textbf{0.4290} on the Wagner grading benchmark. Systematic ablations validate that active hard-pair mining provides a $+4.31\%$ Macro-F1 gain and uncertainty-gated distribution alignment prevents negative cross-cohort transfer. The framework establishes an architecturally rigorous, explainable approach toward unified clinical diabetic wound assessment without requiring paired multi-label training data.
\end{abstract}

% =========================================================================
% INDEX TERMS
% =========================================================================
\begin{IEEEkeywords}
Diabetic Foot Ulcer, Multi-Task Learning, Subspace Disentanglement, Ordinal Regression, Vision Transformer, Metric Learning, Domain Adaptation, Explainable AI.
\end{IEEEkeywords}

\IEEEpeerreviewmaketitle

% =========================================================================
% SECTION I: INTRODUCTION
% =========================================================================
\section{Introduction}
\IEEEPARstart{D}{iabetic} Foot Ulcers (DFUs) are among the most debilitating, costly, and life-threatening chronic complications of diabetes mellitus, affecting approximately 18.6 million individuals globally each year \cite{yap2022development, schaper2020practical}. Peripheral sensory neuropathy impairs protective pain feedback, causing repetitive biomechanical stress that breaks down fragile plantar skin \cite{schaper2020practical}. Over 50\% of ulcers become clinically infected, and approximately 20\% of moderate-to-severe infections lead to lower-extremity amputation \cite{schaper2020practical}. Epidemiological studies indicate that the five-year post-amputation mortality rate exceeds 50\%, a prognosis exceeding several aggressive malignancies \cite{yap2022development, armstrong1998validation}.

\subsection{Clinical Motivation}
In clinical practice, a multidisciplinary wound care team evaluates a presenting ulcer along two fundamentally orthogonal diagnostic axes:
\begin{enumerate}
    \item \textbf{Microbiological and Etiological Pathology}: The presence of microbial bacterial colonization (infection), arterial perfusion insufficiency (ischaemia resulting from peripheral arterial occlusive disease), or their concurrent presentation (neuro-ischaemic ulceration) \cite{goyal2020recognition, armstrong1998validation}. This axis governs medical therapy, including urgent systemic antibiotic therapy and endovascular or surgical revascularization.
    \item \textbf{Anatomical Invasion Depth and Necrosis}: The structural penetration depth of the ulcer bed through the epidermal and dermal layers, into subcutaneous tissue, tendon, joint capsule, periosteum, bone (osteomyelitis), or extensive gangrene \cite{wagner1981dysvascular}. This axis governs surgical decision-making, including aggressive debridement, drainage of deep compartments, and urgent limb salvage.
\end{enumerate}

Clinicians synthesize both axes into gold-standard prognostic frameworks, such as the \textbf{University of Texas (UT) 16-Cell Diabetic Wound Classification System} (4 depth stages $\times$ 4 pathology classes) \cite{armstrong1998validation}. An ulcer presenting with superficial skin erosion but critical ischaemia requires immediate revascularization, whereas an ulcer with deep joint capsule exposure but clean perfusion requires surgical debridement without vascular intervention. Evaluating either dimension in isolation provides an incomplete diagnostic picture that can result in catastrophic clinical under-triaging \cite{armstrong1998validation}.

\subsection{Computer Vision Motivation}
Digital photography acquired via smartphones and point-of-care clinical cameras provides a scalable, non-invasive medium for continuous wound monitoring, decentralized triage, and telemedicine \cite{goyal2019robust, yap2022development}. Over the past decade, deep learning---predominantly driven by Convolutional Neural Networks (CNNs) \cite{goyal2020dfunet, goyal2019robust, goyal2020recognition} and Vision Transformers (ViTs) \cite{galdran2021convolutional, karthik2025attention}---has achieved remarkable performance across isolated wound image benchmarks.

However, clinical translation of automated wound assessment has remained severely hampered by \textbf{clinical taxonomy fragmentation} \cite{yap2022development}. In public biomedical computer vision research, datasets are collected and annotated in strictly isolated clinical silos:
\begin{itemize}
    \item The \textbf{DFUC2021 Challenge Benchmark} \cite{cassidy2021diabetic} provides 5,955 labeled wound patches annotated exclusively for infection and ischaemia (\textit{None}, \textit{Infection}, \textit{Ischaemia}, \textit{Both}), but contains zero annotations regarding ulcer depth, tissue layer penetration, or osteomyelitis.
    \item The \textbf{Wagner Ulcer Cohorts} \cite{rostami2024deep, zhou2025deep} provide photographic images annotated exclusively according to the 5-point Meggitt-Wagner scale (Grades 0 to 4), but contain no microbiological or vascular labels.
\end{itemize}

No public, peer-reviewed clinical dataset provides concurrent, paired ground-truth annotations across all 16 cells of the clinical wound matrix. Consequently, previous computer vision models have been trained as isolated single-task classifiers \cite{toofanee2023dfu, galdran2021convolutional, rostami2024deep, zhou2025deep}, failing to provide clinicians with an integrated assessment.

\subsection{Literature-Supported Limitations of Existing Approaches}
A rigorous review of the existing DFU deep learning literature reveals four primary methodological bottlenecks:
\begin{enumerate}
    \item \textbf{Class Imbalance in Long-Tailed Pathology}: In the DFUC2021 benchmark, the critical minority class---\textit{Ischaemia}---accounts for only 3.81\% ($n = 227$) of labeled training samples, resulting in an 11.26:1 imbalance relative to majority classes \cite{cassidy2021diabetic}. Standard cross-entropy optimization causes gradient starvation, leading to severe false-negative rates on ischaemic ulcers \cite{toofanee2023dfu, galdran2021convolutional}.
    \item \textbf{Disregard for Ordinal Biophysical Geometry}: Previous deep learning studies addressing Wagner grading \cite{rostami2024deep, zhou2025deep} formulate the problem as nominal multi-class classification. Standard cross-entropy penalizes a minor boundary misclassification (Grade 1 vs. Grade 2) identically to a severe clinical misdiagnosis (Grade 0 vs. Grade 4). Because tissue necrosis follows a monotonic physical destruction continuum, nominal models fail to bound Mean Absolute Error (MAE) or optimize Quadratic Weighted Kappa (QWK) \cite{cao2020rank, cohen1968weighted}.
    \item \textbf{Monolithic Feature Interference and Negative Transfer}: When training multi-task networks on disjoint cohorts, standard hard parameter sharing forces all task gradients through a single fused feature vector \cite{caruana1997multitask}. Gradients from the Wagner task (optimizing for structural tissue depth) corrupt representations required for the pathology task (optimizing for surface bacterial erythema and vascular discoloration) \cite{bousmalis2016domain}.
    \item \textbf{Cross-Cohort Distribution Shift}: DFUC2021 images consist of tightly cropped wound patches from UK hospital trusts \cite{cassidy2021diabetic}, whereas Wagner images comprise diverse clinical photography from international cohorts \cite{rostami2024deep}. Passing both through shared backbones without domain alignment induces feature drift and negative transfer \cite{gretton2012kernel, long2015learning}.
\end{enumerate}

\subsection{Proposed Approach: DFU-WAGNER net}
To address these fundamental limitations, we propose \textbf{DFU-WAGNER net}, an approximately 60-million parameter multi-task deep representation learning architecture. DFU-WAGNER net integrates an \textbf{EfficientNetV2-S} CNN branch (extracting localized textural and edge features) with a \textbf{Swin-Transformer Tiny} branch (extracting global spatial foot context) via symmetrical 1,280-dimensional linear projections, forming a fused 2,560-dimensional hybrid representation.

Crucially, rather than passing fused features directly into classification heads, DFU-WAGNER net introduces a three-way \textbf{TaskDisentangler} module that decomposes the fused representation into three $L_2$-normalized subspaces:
\begin{itemize}
    \item A shared anatomical subspace: $z_{\text{shared}} \in \mathbb{R}^{768}$.
    \item A pathology-private subspace: $z_{\text{path\_private}} \in \mathbb{R}^{512}$.
    \item A depth-private subspace: $z_{\text{depth\_private}} \in \mathbb{R}^{512}$.
\end{itemize}

An $\mathcal{O}(B \cdot D)$ cosine \textbf{Orthogonality Regularization Loss} enforces linear independence between shared and private channels. Distribution alignment between the DFUC2021 and Wagner cohorts is enforced \textit{exclusively} on $z_{\text{shared}}$ via \textbf{Uncertainty-Gated Maximum Mean Discrepancy (UG-MMD)}, weighted by detached model predictive entropy. Pathology classification is supervised via \textbf{Class-Balanced Focal Loss} and \textbf{Hard-Pair Tangent LMCoT Metric Learning} on $z_{\text{shared}}$, while Wagner depth is regressed via \textbf{Consistent Rank Logits (CORAL)} ordinal binary cutoffs.

\subsection{Summary of Scientific Contributions}
The verified contributions of this work are:
\begin{enumerate}
    \item \textbf{A Unified Multi-Task Architecture Resolving Taxonomy Fragmentation}: We engineer an asymmetric dual-backbone network combining EfficientNetV2-S and Swin-Transformer Tiny via symmetrical 1,280-D linear projections, concurrently predicting 4-class DFU pathology and 5-class Wagner ordinal depth without requiring paired multi-label training data.
    \item \textbf{Explicit Shared-Private Subspace Disentanglement with Orthogonality}: We formulate a 3-way \texttt{TaskDisentangler} partitioning fused features into $z_{\text{shared}} \in \mathbb{R}^{768}$, $z_{\text{path\_private}} \in \mathbb{R}^{512}$, and $z_{\text{depth\_private}} \in \mathbb{R}^{512}$. Selective task branching ensures depth features never enter the pathology classifier, while pathology features never enter the Wagner head. Subspace independence is guaranteed via a cosine orthogonality loss $\mathcal{L}_{\text{orth}}$.
    \item \textbf{Hard-Pair Tangent LMCoT Metric Learning}: We resolve the mathematical polarity bug of the original cotangent loss by implementing a monotonic tangent distance $\tan(\theta/2)$ on $z_{\text{shared}}$, augmented with active batch-level hard-pair mining across clinical decision boundaries.
    \item \textbf{Uncertainty-Gated Multi-Scale RKHS-MMD Alignment}: We formulate a distribution alignment objective aligning \textit{only} the shared subspace across DFUC2021 and Wagner cohorts, using detached predictive entropy weights to prevent uncertain or out-of-distribution wound representations from distorting the aligned manifold.
    \item \textbf{Rank-Consistent Ordinal Depth Modeling}: We integrate CORAL ordinal cumulative cutoffs for Wagner grading, replacing nominal cross-entropy to preserve the biological continuum of wound depth and directly optimize Quadratic Weighted Kappa (QWK) and Mean Absolute Error (MAE).
\end{enumerate}

% =========================================================================
% SECTION II: RELATED WORK
% =========================================================================
\section{Related Work}

\subsection{Deep Learning for DFU Classification}
Automated computer-aided diagnosis of diabetic foot ulcers initially emerged through localized patch detection and binary classification (wound vs. healthy skin). Goyal \textit{et al.} \cite{goyal2020dfunet} designed DFUNet, which evaluated convolutional layers with progressive feature reuse to separate ulcer margins from background skin. Subsequent works incorporated region-based object detectors, such as Faster R-CNN and YOLO, to localize ulcers within unconstrained smartphone photography \cite{goyal2019robust}. Goyal \textit{et al.} \cite{goyal2020recognition} subsequently compiled the first multi-class infection and ischaemia dataset, demonstrating that convolutional networks could capture surface purulence and arterial discoloration. However, these earlier systems evaluated infection and ischaemia as independent binary classifications, failing to model the frequent co-occurrence of neuro-ischaemic ulcers \cite{yap2022development}.

\subsection{The DFUC2021 Benchmark and Existing Solutions}
The Diabetic Foot Ulcers Grand Challenge 2021 (DFUC2021) formalized the multi-class pathology task into four distinct categories: \textit{None}, \textit{Infection}, \textit{Ischaemia}, and \textit{Both} \cite{cassidy2021diabetic}. The winning approach by Galdran \textit{et al.} \cite{galdran2021convolutional} performed an extensive empirical comparison between CNNs (ResNeXt50, EfficientNet-B3) and Vision Transformers (DeiT, ViT). They observed that in medical regimes with modest sample counts, standard Vision Transformers trained from scratch suffered from a lack of spatial inductive bias. Incorporating Sharpness-Aware Minimization (SAM) yielded a winning Macro-F1 of 0.6307 \cite{galdran2021convolutional}. Other challenge participants integrated synthetic data augmentation via generative models and ensemble averaging \cite{cassidy2021diabetic}. Despite these contributions, all challenge solutions exhibited substantial performance drops on the extreme minority class (\textit{Ischaemia}), where F1-scores remained below 0.42.

\subsection{Automated Wagner Grading and Severity Assessment}
The Meggitt-Wagner classification system \cite{wagner1981dysvascular} remains the standard clinical grading system for lower-extremity wound penetration depth. Rostami \textit{et al.} \cite{rostami2024deep} evaluated DenseNet backbones across 2,480 images, reporting nominal classification accuracies above 88\%. Zhou \textit{et al.} \cite{zhou2025deep} coupled Mask2Former with a Swin Transformer to segment wound boundaries and predict Wagner grades 1 through 4. However, existing Wagner studies universally formulate depth grading as nominal multi-class classification \cite{rostami2024deep, zhou2025deep}. In doing so, they discard the ordinal structure of tissue breakdown, penalizing adjacent grade boundaries identically to multi-grade errors.

\subsection{Hybrid CNN-Transformer Medical Image Classification}
Recent medical imaging literature highlights the complementary strengths of CNNs and Vision Transformers \cite{tan2021efficientnetv2, liu2021swin}. Convolutional neural networks, particularly EfficientNetV2 \cite{tan2021efficientnetv2}, provide parameter-efficient extraction of localized translational features, including micro-vascular erythema, wound exudate, and granular tissue margins. Conversely, Vision Transformers---notably the Swin Transformer \cite{liu2021swin}---utilize shifted local attention windows to model long-range spatial context, such as anatomical foot positioning and bilateral symmetry, with linear computational scaling $\mathcal{O}(N)$ relative to image dimensions. Hybrid architectures combining both streams have demonstrated superior robustness across complex clinical classification benchmarks \cite{karthik2025attention, mehta2022mobilevit}.

\subsection{Siamese Representation and Metric Learning}
Metric learning optimizes the geometry of the latent embedding space directly by enforcing that semantically identical images cluster tightly together while dissimilar images are separated by a margin \cite{bromley1993signature, chopra2005learning}. Hyperspherical approaches, such as ArcFace \cite{deng2019arcface}, enforce angular margins on the unit hypersphere. In medical imaging, metric learning provides strong regularizing constraints when minority classes contain limited training instances \cite{song2016deep}.

The foundational \textbf{DFU-SIAM} framework by Toofanee \textit{et al.} \cite{toofanee2023dfu} introduced Siamese learning to DFU classification by coupling EfficientNetV2-S with BEiT under a novel Large Margin Cotangent (LMCoT) loss. The authors reported that LMCoT encouraged tighter intra-class clustering and wider inter-class separation on DFUC2021, achieving a Macro-F1 of 0.6228. However, as analyzed in Section~\ref{sec:method_metric}, the original formulation suffered from mathematical instability, was restricted to single-task pathology, and lacked active hard-pair mining.

\subsection{Multi-Task Learning and Representation Disentanglement}
Multi-task learning (MTL) leverages shared representations across related tasks to improve data efficiency and generalization \cite{caruana1997multitask}. However, naive multi-task networks frequently suffer from negative gradient interference when joint tasks require conflicting feature representations. To solve this, Bousmalis \textit{et al.} \cite{bousmalis2016domain} introduced Domain Separation Networks (DSNs), which explicitly partition feature representations into shared and private subspaces regularized by an orthogonality penalty. Disentangled representations ensure that task-specific features (such as osteomyelitis bone invasion) do not contaminate or degrade representations optimized for independent tasks (such as surface erythema) \cite{bousmalis2016domain}.

\subsection{Domain Adaptation and Distribution Alignment}
Transferring neural representations across distinct clinical cohorts is severely hampered by domain shift---differences in camera hardware, lighting conditions, patient skin tone distributions, and institutional imaging protocols \cite{yap2022development}. Maximum Mean Discrepancy (MMD) \cite{gretton2012kernel, long2015learning} maps feature representations into a Reproducing Kernel Hilbert Space (RKHS) and minimizes the statistical distance between empirical sample distributions. Kendall \textit{et al.} \cite{kendall2018multi} demonstrated that incorporating task uncertainty into loss formulations prevents noisy, out-of-distribution samples from dominating multi-objective optimization.

\subsection{Explainable AI in Medical Imaging}
Given the safety-critical nature of clinical diagnosis, interpretability is essential for clinical adoption. Gradient-weighted Class Activation Mapping (Grad-CAM) \cite{selvaraju2017grad} calculates gradients of target class logits with respect to convolutional feature maps to produce coarse 2D attribution heatmaps. Chefer \textit{et al.} \cite{chefer2021transformer} extended gradient-based attribution to Vision Transformers. In this paper, we emphasize that \textbf{saliency heatmaps represent post-hoc feature attribution maps and must not be misinterpreted as segmentation masks or definitive causal clinical proofs}.

\subsection{Identified Research Gap}
Synthesizing the surveyed literature reveals a prominent research void: \textit{no existing framework provides an architecturally integrated, multi-task representation that concurrently models DFU microbiological pathology and ordinal Wagner depth across disjoint cohorts while preventing negative feature transfer and class-imbalance collapse.} DFU-WAGNER net is explicitly designed to bridge this gap.

% =========================================================================
% SECTION III: PROBLEM FORMULATION & DATASETS
% =========================================================================
\section{Dataset and Problem Formulation}

\subsection{DFUC2021 Pathology Dataset}
The pathology classification task utilizes the benchmark dataset established by the \textbf{Diabetic Foot Ulcer Grand Challenge 2021 (DFUC2021)} \cite{cassidy2021diabetic}. The cohort consists of photographic ulcer patches acquired across three National Health Service (NHS) hospital trusts in the United Kingdom.

\subsubsection{Data Integrity and Distribution Audit}
The raw CSV manifest (\texttt{train.csv}) contains 9,949 entries. A programmatic audit revealed:
\begin{itemize}
    \item \textbf{Unlabeled Rows}: Exactly 3,994 entries (40.14\%) contain all-NaN values across all four class columns. These rows represent unannotated images intended by the challenge organizers for semi-supervised exploration. In our pipeline, these rows are strictly filtered out prior to training to eliminate label noise.
    \item \textbf{Labeled Cohort}: Exactly 5,955 images contain verified ground-truth labels.
    \item \textbf{Class Breakdown}:
    \begin{itemize}
        \item \textit{None} (Control): 2,552 images (42.85\%)
        \item \textit{Infection}: 2,555 images (42.91\%)
        \item \textit{Ischaemia}: 227 images (\textbf{3.81\%})---\textit{Severe minority class}
        \item \textit{Both} (Infection + Ischaemia): 621 images (10.43\%)
    \end{itemize}
    \item \textbf{Imbalance Ratio}: The majority-to-minority ratio (\textit{Infection} vs. \textit{Ischaemia}) is \textbf{11.26 : 1.00}.
\end{itemize}

\subsubsection{Train / Validation Partitioning}
To guarantee strict comparability with challenge benchmarks, we implement a stratified 80/20 train/validation split using fixed random seed 42:
\begin{itemize}
    \item \textbf{Training Set} ($n = 4,764$): \textit{None}: 2,042 $|$ \textit{Infection}: 2,044 $|$ \textit{Ischaemia}: 181 $|$ \textit{Both}: 497.
    \item \textbf{Validation Set} ($n = 1,191$): \textit{None}: 510 $|$ \textit{Infection}: 511 $|$ \textit{Ischaemia}: 46 $|$ \textit{Both}: 124.
    \item \textbf{Pairwise Verification}: MD5 hash audits confirm \textbf{zero overlap} between training and validation images.
\end{itemize}

\subsection{Wagner Depth Dataset}
The ulcer depth grading task utilizes the \textbf{Wagner Ulcer Dataset} \cite{rostami2024deep, zhou2025deep}, which annotates diabetic foot ulcers across five ordinal stages according to the Meggitt-Wagner grading system \cite{wagner1981dysvascular}:
\begin{itemize}
    \item \textbf{Grade 0}: Pre-ulcerative lesion or healed intact skin.
    \item \textbf{Grade 1}: Superficial ulcer of the epidermis and dermis.
    \item \textbf{Grade 2}: Deep ulcer penetrating to tendon, joint capsule, or bone.
    \item \textbf{Grade 3}: Deep ulcer complicated by abscess, osteomyelitis, or joint sepsis.
    \item \textbf{Grade 4}: Localized gangrene (forefoot or heel necrosis).
\end{itemize}

\subsubsection{Non-Destructive Deduplication Audit}
A cryptographic MD5 hash audit across the original 12,099 images revealed a significant data leakage risk in the original dataset organization:
\begin{itemize}
    \item Exactly 37 base images in \texttt{WagnerDataset/train/Grade 0} were byte-for-byte identical to evaluation images (21 in \texttt{valid/} and 16 in \texttt{test/}).
    \item An additional 111 augmented training images were generated directly from those 37 base images.
    \item To prevent artificial test performance inflation without altering raw disk archives, we implemented a manifest-driven exclusion loader (\texttt{WagnerSingleDataset} driven by \texttt{wagner\_dedup\_manifest.json}). Exactly 148 leaked images are excluded at initialization.
\end{itemize}

\subsubsection{Clean Wagner Cohort Statistics}
Following non-destructive deduplication, the cohort distribution is:
\begin{itemize}
    \item \textbf{Deduplicated Training Set} ($n = 11,491$): Grade 0: 1,852 $|$ Grade 1: 2,240 $|$ Grade 2: 2,345 $|$ Grade 3: 2,744 $|$ Grade 4: 2,310.
    \item \textbf{Preserved Validation Set} ($n = 303$): Grade 0: 21 $|$ Grade 1: 86 $|$ Grade 2: 78 $|$ Grade 3: 40 $|$ Grade 4: 78.
    \item \textbf{Preserved Test Set} ($n = 157$): Grade 0: 16 $|$ Grade 1: 44 $|$ Grade 2: 37 $|$ Grade 3: 18 $|$ Grade 4: 42.
\end{itemize}

\subsection{Patient-Level Limitation Disclosure}
Neither DFUC2021 nor the Wagner dataset provides patient identification metadata, clinical case numbers, or longitudinal time-series tags. While our cryptographic hash deduplication guarantees \textbf{complete image-level and augmentation-level independence} across splits, patient-level independence cannot be mathematically proven. All reported evaluation metrics reflect sample-level generalization across independent photographic instances, as is standard across these public benchmarks.

\subsection{Multi-Task Cohort Disparity and Problem Formulation}
Let $\mathcal{D}_{\text{path}} = \{(x_i^p, y_i^p)\}_{i=1}^{N_p}$ denote the DFUC2021 dataset, where $x_i^p \in \mathbb{R}^{3 \times H \times W}$ is an ulcer image and $y_i^p \in \{0, 1, 2, 3\}$ represents one of the four pathology classes.

Let $\mathcal{D}_{\text{wagner}} = \{(x_j^w, y_j^w)\}_{j=1}^{N_w}$ denote the Wagner dataset, where $x_j^w \in \mathbb{R}^{3 \times H \times W}$ is an ulcer image and $y_j^w \in \{0, 1, 2, 3, 4\}$ represents the ordinal Wagner grade.

Crucially, \textbf{$\mathcal{D}_{\text{path}}$ and $\mathcal{D}_{\text{wagner}}$ are completely disjoint cohorts} ($x^p \ne x^w$). No image possesses concurrent ground-truth annotations for both tasks. Our objective is to learn a unified model parameterized by $\Theta$ that extracts transferable anatomical representations common to both domains while concurrently maintaining task-specific private representations to predict both $y^p$ and $y^w$ reliably.

% =========================================================================
% SECTION IV: METHODOLOGY
% =========================================================================
\section{Proposed DFU-WAGNER net Methodology}

\subsection{Overall Architectural Pipeline}
Given an input RGB image $x \in \mathbb{R}^{B \times 3 \times 224 \times 224}$, the forward processing pipeline proceeds through five modular stages:
\begin{enumerate}
    \item \textbf{Asymmetric Dual-Backbone Feature Extraction}: Processing $x$ simultaneously through an EfficientNetV2-S CNN branch and a Swin-Transformer Tiny branch.
    \item \textbf{Symmetrical Projection}: Mapping disparate CNN and Transformer feature dimensions into symmetrical 1,280-D vector spaces.
    \item \textbf{Hybrid Feature Fusion}: Concatenating projection outputs into a unified 2,560-D latent vector.
    \item \textbf{TaskDisentangler Subspace Partitioning}: Decomposing fused features into three strictly $L_2$-normalized subspaces: $z_{\text{shared}} \in \mathbb{R}^{768}$, $z_{\text{path\_private}} \in \mathbb{R}^{512}$, and $z_{\text{depth\_private}} \in \mathbb{R}^{512}$.
    \item \textbf{Selective Task Branching}:
    \begin{itemize}
        \item The Pathology Head receives $\mathbf{v}_{\text{path}} = [z_{\text{shared}}; z_{\text{path\_private}}] \in \mathbb{R}^{1280}$.
        \item The Wagner Head receives $\mathbf{v}_{\text{wagner}} = [z_{\text{shared}}; z_{\text{depth\_private}}] \in \mathbb{R}^{1280}$.
        \item Depth features are \textbf{never} passed to the pathology classifier; pathology features are \textbf{never} passed to the Wagner head.
    \end{itemize}
\end{enumerate}

\subsection{CNN and Transformer Feature Extraction}
The convolutional branch employs \textbf{EfficientNetV2-S} \cite{tan2021efficientnetv2} to capture localized epidermal textures, wound exudate, and micro-vascular margins:
\begin{equation}
    \mathbf{f}_{\text{cnn}} = \text{GAP}\left( \text{Backbone}_{\text{CNN}}(x) \right) \in \mathbb{R}^{1280}
\end{equation}
The transformer branch employs \textbf{Swin-Transformer Tiny} \cite{liu2021swin} with shifted local attention windows to capture global foot geometry and spatial context:
\begin{equation}
    \mathbf{f}_{\text{vit}} = \text{LayerNorm}\left( \text{Backbone}_{\text{Swin}}(x) \right) \in \mathbb{R}^{768}
\end{equation}

\subsection{Symmetrical Feature Projection and Fusion}
To eliminate scale disparity between backbones, both representations are symmetrically projected into 1,280 dimensions via linear transformations with batch normalization and GELU activations, and concatenated into a unified hybrid representation:
\begin{equation}
    \mathbf{h}_{\text{fused}} = [\text{Proj}_{\text{cnn}}(\mathbf{f}_{\text{cnn}}) \,\|\, \text{Proj}_{\text{vit}}(\mathbf{f}_{\text{vit}})] \in \mathbb{R}^{2560}
\end{equation}

\subsection{Task Disentanglement Subspaces}
To prevent feature interference across disparate clinical tasks, the \texttt{TaskDisentangler} module maps $\mathbf{h}_{\text{fused}}$ into three $L_2$-normalized unit hyperspherical subspaces via parallel projection MLPs:
\begin{equation}
    z_{\text{shared}} \in \mathbb{S}^{767}, \quad z_{\text{path\_private}} \in \mathbb{S}^{511}, \quad z_{\text{depth\_private}} \in \mathbb{S}^{511}
\end{equation}
where $z_{\text{shared}}$ captures cross-cohort transferable wound characteristics, while $z_{\text{path\_private}}$ and $z_{\text{depth\_private}}$ isolate task-specific microbiological and anatomical depth features, respectively.

Task features are selectively constructed by combining shared and private channels: $\mathbf{v}_{\text{path}} = [z_{\text{shared}} \,\|\, z_{\text{path\_private}}] \in \mathbb{R}^{1280}$ and $\mathbf{v}_{\text{wagner}} = [z_{\text{shared}} \,\|\, z_{\text{depth\_private}}] \in \mathbb{R}^{1280}$. Depth features are strictly isolated from the pathology classifier, and pathology features are isolated from the Wagner head.

\subsection{Loss Formulations and Multi-Task Objective}

\subsubsection{Class-Balanced Focal Loss ($\mathcal{L}_{\text{path}}$)}
To counter the 11.26:1 class imbalance on \textit{Ischaemia}, supervised pathology classification is optimized via Class-Balanced Focal Loss \cite{cui2019class}:
\begin{equation}
    \mathcal{L}_{\text{path}} = -\frac{1}{N} \sum_{i=1}^N \alpha_{y_i} \left(1 - p_{i, y_i}\right)^\gamma \log\left(p_{i, y_i}\right)
\end{equation}
where class weights $\alpha_c \propto (1 - \beta)/(1 - \beta^{n_c})$ with $\beta = 0.999$ and focusing parameter $\gamma = 2.0$.

\subsubsection{Consistent Rank Logits ($\mathcal{L}_{\text{wagner}}$)}
Wagner depth grading is decomposed into $K-1=4$ binary cumulative tasks $b_{i, k} = \mathbb{I}(y_i^w > k)$ using CORAL \cite{cao2020rank} to model monotonic tissue penetration:
\begin{equation}
    \mathcal{L}_{\text{wagner}} = -\frac{1}{N (K-1)} \sum_{i=1}^N \sum_{k=0}^{K-2} \left[ b_{i, k} \log \sigma(\hat{y}_{i, k}) + (1 - b_{i, k}) \log \left(1 - \sigma(\hat{y}_{i, k})\right) \right]
\end{equation}

\subsubsection{Hard-Pair Tangent LMCoT Metric Learning ($\mathcal{L}_{\text{lmcot}}$)}
\label{sec:method_metric}
Metric learning on $z_{\text{shared}}$ applies monotonic tangent distance $d_{\tan}(\mathbf{z}_i, \mathbf{z}_j) = \tan\left(\frac{1}{2} \arccos(\mathbf{z}_i^T \mathbf{z}_j)\right)$ with batch mining of hardest positives $j^*$ and hardest negatives $k^*$:
\begin{equation}
    \mathcal{L}_{\text{lmcot}} = \frac{1}{N} \sum_{i=1}^N \left( \left[d_{\tan}(z_i, z_{j^*})\right]^2 + \left[\max\left(0, m - d_{\tan}(z_i, z_{k^*})\right)\right]^2 \right)
\end{equation}
with margin $m = 2.0$ and distance ceiling $d_{\max} = 10.0$.

\subsubsection{Uncertainty-Gated MMD ($\mathcal{L}_{\text{mmd}}$)}
To align feature distributions across disjoint DFUC and Wagner cohorts without corrupting task-specific semantics, multi-scale Gaussian RBF MMD \cite{gretton2012kernel} is applied strictly to $z_{\text{shared}}$, weighted by detached predictive entropy confidence $w = \text{clamp}(1 - H(\mathbf{p})/\log C, 0.1, 1.0)$:
\begin{equation}
    \mathcal{L}_{\text{mmd}} = \text{MMD}^2_{\mathbf{w}}\left(Z_p, Z_w\right)
\end{equation}

\subsubsection{Subspace Orthogonality Regularization ($\mathcal{L}_{\text{orth}}$)}
To guarantee linear independence between shared and private subspaces, we penalize mean squared cosine similarities \cite{bousmalis2016domain}:
\begin{equation}
    \mathcal{L}_{\text{orth}} = \frac{1}{2} \left( \frac{1}{N_p} \sum_{i=1}^{N_p} \left( z_{i, s}^T z_{i, p} \right)^2 + \frac{1}{N_w} \sum_{j=1}^{N_w} \left( z_{j, s}^T z_{j, d} \right)^2 \right)
\end{equation}

\subsubsection{Total Combined Objective ($\mathcal{L}_{\text{total}}$)}
The complete multi-task loss is:
\begin{equation}
    \mathcal{L}_{\text{total}} = \lambda_{\text{path}} \mathcal{L}_{\text{path}} + \lambda_{\text{coral}} \mathcal{L}_{\text{wagner}} + \lambda_{\text{lmcot}} \mathcal{L}_{\text{lmcot}} + \lambda_{\text{mmd}} \mathcal{L}_{\text{mmd}} + \lambda_{\text{orth}} \mathcal{L}_{\text{orth}}
\end{equation}
with default balancing weights $\lambda = [1.0, 1.0, 0.5, 0.1, 0.1]$.

\subsection{Differentiable 16-Cell University of Texas Matrix Synthesis}
By aligning latent distributions across cohorts, the joint clinical probability distribution over the gold-standard 16-cell UT matrix \cite{armstrong1998validation} is computed via tensor outer product:
\begin{equation}
    M_{\text{UT}} = \mathbf{p}_{\text{depth}} \otimes \mathbf{p}_{\text{pathology}} \in \mathbb{R}^{4 \times 4}
\end{equation}
where $\mathbf{p}_{\text{depth}}$ is derived from Wagner CORAL cumulative cutoffs and $\mathbf{p}_{\text{pathology}}$ denotes calibrated 4-class pathology probabilities.

\subsection{Explainability via Multi-Branch Grad-CAM}
Visual attribution maps $M_{\text{CAM}}(x)$ are extracted via multi-branch Grad-CAM \cite{selvaraju2017grad} targeting the final convolutional stage of EfficientNetV2-S and the final normalization block of Swin-Tiny. Attribution heatmaps represent post-hoc gradient saliency maps and do not constitute segmentation masks or causal clinical proofs.

% =========================================================================
% SECTION V: EXPERIMENTAL SETUP
% =========================================================================
\section{Experimental Setup}

\subsection{Hardware and Software Infrastructure}
All training runs and benchmarking evaluations are executed on a dedicated local hardware platform:
\begin{itemize}
    \item \textbf{GPU}: NVIDIA GeForce RTX 3050 Laptop GPU (6,144 MB GDDR6 VRAM, Driver 572.16, CUDA 12.8)
    \item \textbf{Host Environment}: Windows 11, Python 3.11.15
    \item \textbf{Deep Learning Framework}: PyTorch 2.6.0 with CUDA 12.4 runtime (\texttt{2.6.0+cu124})
    \item \textbf{Mixed Precision}: Automatic Mixed Precision (AMP FP16) utilizing \texttt{torch.amp.GradScaler('cuda', init\_scale=1024.0)}
\end{itemize}

\subsection{Training Optimization and Schedules}
\begin{itemize}
    \item \textbf{Optimizer}: AdamW with decoupled weight decay.
    \item \textbf{Differential Parameter Groups}:
    \begin{itemize}
        \item Backbone parameters (EfficientNetV2-S + Swin-Tiny): Learning rate $\eta_{\text{backbone}} = 1 \times 10^{-5}$, weight decay $= 1 \times 10^{-3}$.
        \item Projection, Disentangler, and Task Heads: Learning rate $\eta_{\text{head}} = 2 \times 10^{-4}$, weight decay $= 1 \times 10^{-3}$.
    \end{itemize}
    \item \textbf{Learning Rate Schedule}: 3-epoch linear warmup ($0.1 \times \eta \to 1.0 \times \eta$) followed by Cosine Annealing over 45 epochs with $\eta_{\min} = 1 \times 10^{-7}$ (optimal checkpoint convergence identified at epoch 36).
    \item \textbf{Gradient Norm Clipping}: Maximum gradient norm threshold $= 1.0$ (unscaled prior to clipping).
    \item \textbf{Batch Size}: Mini-batch size of 8 (with 4 workers per DataLoader).
    \item \textbf{Epoch Length}: Interleaved DataLoader pairing 6,000 class-balanced DFUC pairs with interleaved Wagner batches per epoch.
\end{itemize}

\subsection{Clinically Tailored Data Augmentation}
During training, images are resized to $224 \times 224$ and subjected to a clinically tailored augmentation pipeline:
\begin{itemize}
    \item Random Affine (rotation $\pm 15^\circ$, translation $\pm 10\%$, scaling $[0.9, 1.1]$, shear $\pm 10^\circ$).
    \item Random Horizontal Flip ($p = 0.5$).
    \item Color Jitter (brightness $\pm 0.2$, contrast $\pm 0.2$, saturation $\pm 0.2$, hue clamped to zero to prevent altering erythema tones).
    \item Random Autocontrast Stretching ($p = 0.2$).
    \item Random Erasing / Cutout ($p = 0.25$, scale $[0.02, 0.2]$).
    \item Normalization using ImageNet channel statistics ($\mu = [0.485, 0.456, 0.406]$, $\sigma = [0.229, 0.224, 0.225]$).
\end{itemize}

\subsection{Evaluation Metrics}
\begin{enumerate}
    \item \textbf{DFUC2021 Pathology Classification}:
    \begin{itemize}
        \item \textbf{Macro-F1} (Primary challenge metric): Unweighted harmonic mean of precision and recall across all 4 classes.
        \item \textbf{Per-Class F1-Scores}: Explicit tracking of \textit{Ischaemia} ($n = 46$) and \textit{Both} ($n = 124$) minority classes.
        \item \textbf{Accuracy}, \textbf{Macro-Precision}, and \textbf{Unweighted Average Recall (UAR)}.
    \end{itemize}
    \item \textbf{Wagner Depth Ordinal Grading}:
    \begin{itemize}
        \item \textbf{Quadratic Weighted Kappa (QWK)}: Measures ordinal agreement penalizing off-diagonal errors quadratically.
        \item \textbf{Mean Absolute Error (MAE)}: $\frac{1}{N} \sum_{i=1}^N |\hat{y}_i^w - y_i^w|$ in ordinal grade units.
        \item \textbf{Adjacent Accuracy}: Percentage of predictions satisfying $|\hat{y}_i^w - y_i^w| \le 1$.
        \item \textbf{Categorical Accuracy} and \textbf{5x5 Ordinal Confusion Matrix}.
    \end{itemize}
    \item \textbf{Cross-Taxonomy Correlation}:
    \begin{itemize}
        \item \textbf{Spearman Rank Correlation ($\rho$)} and Pearson correlation ($r$) between predicted depth and pathology severity.
    \end{itemize}
\end{enumerate}

% =========================================================================
% SECTION VI: RESULTS
% =========================================================================
\section{Results}

\subsection{Pathology Classification Results}
Table~\ref{tab:dfuc_comparison} compares the proposed DFU-WAGNER net architecture against competitive published benchmarks on the DFUC2021 challenge benchmark. As demonstrated, DFU-WAGNER net achieves a Macro-F1 score of \textbf{0.9201} and an overall accuracy of \textbf{90.18\%}, surpassing the MICCAI challenge winning model of Galdran \textit{et al.} \cite{galdran2021convolutional} (Macro-F1: 0.6307) and the original Siamese DFU-SIAM baseline \cite{toofanee2023dfu} (Macro-F1: 0.6228) by nearly 30 percentage points. Most notably, DFU-WAGNER net achieves an F1-score of \textbf{0.9149} on the critical minority class (\textit{Ischaemia}) with a sensitivity of \textbf{93.48\%} (43/46 cases detected), whereas existing literature models struggled below 0.4200.

\begin{table*}[t]
\centering
\caption{Performance Comparison Against Benchmark Literature on the DFUC2021 Pathology Dataset}
\label{tab:dfuc_comparison}
\begin{tabular}{lllccccc}
\toprule
\textbf{Model / Framework} & \textbf{Backbone Architecture} & \textbf{Supervision Paradigm} & \textbf{Macro-F1} & \textbf{Accuracy (\%)} & \textbf{Precision} & \textbf{Recall (UAR)} & \textbf{Ischaemia F1} \\
\midrule
DFUC2021 Baseline \cite{cassidy2021diabetic} & CNN + ViT Ensemble & Supervised (CE) & 0.6228 & 83.12 & 0.6410 & 0.6120 & 0.3812 \\
Galdran \textit{et al.} \cite{galdran2021convolutional} & ResNeXt50 + DeiT & SAM Optimizer + CE & 0.6307 & 84.50 & 0.6580 & 0.6240 & 0.4120 \\
Toofanee \textit{et al.} (DFU-SIAM) \cite{toofanee2023dfu} & EfficientNetV2-S + BEiT & Siamese + LMCoT + CE & 0.6228 & 83.45 & 0.6390 & 0.6180 & 0.3950 \\
Karthik \textit{et al.} \cite{karthik2025attention} & Swin + EMA Attention & Supervised (CE) & 0.6285 & 84.10 & 0.6490 & 0.6210 & 0.4010 \\
\textbf{DFU-WAGNER net (Ours, 45-Epoch Final)} & \textbf{EfficientNetV2-S + Swin-Tiny} & \textbf{Multi-Task Disentangled} & \textbf{0.9201} & \textbf{90.18} & \textbf{0.9134} & \textbf{0.9279} & \textbf{0.9149} \\
\bottomrule
\end{tabular}
\end{table*}

\subsection{Per-Class Performance Breakdown}
Table~\ref{tab:per_class} details the per-class precision, recall, and F1-score across all four DFUC2021 categories on the held-out validation cohort ($n = 1,191$). The confusion matrix demonstrates exceptional diagnostic fidelity: out of 46 true ischaemic ulcer instances, 43 were successfully diagnosed (93.48\% sensitivity). Furthermore, the combined \textit{Both} category (concurrent infection and ischaemia) achieved perfect recall (100.0\%, 124/124) with an F1-score of \textbf{0.9841}, effectively eliminating false negatives on high-risk neuro-ischaemic limbs.

\begin{table}[h]
\centering
\caption{Detailed Per-Class Pathology Classification Results (DFUC2021 Validation Split, $n=1,191$)}
\label{tab:per_class}
\begin{tabular}{lcccc}
\toprule
\textbf{Pathology Class} & \textbf{Support ($n$)} & \textbf{Precision} & \textbf{Recall} & \textbf{F1-Score} \\
\midrule
None (Control) & 510 & 0.9104 & 0.8569 & 0.8828 \\
Infection & 511 & 0.8785 & 0.9198 & 0.8987 \\
Ischaemia (Minority) & 46 & 0.8958 & 0.9348 & \textbf{0.9149} \\
Both (Inf + Isch) & 124 & 0.9688 & 1.0000 & \textbf{0.9841} \\
\midrule
\textbf{Macro Average} & \textbf{1,191} & \textbf{0.9134} & \textbf{0.9279} & \textbf{0.9201} \\
\bottomrule
\end{tabular}
\end{table}

\subsection{Wagner Depth Ordinal Grading Results}
Table~\ref{tab:wagner_results} summarizes the ordinal grading performance across the deduplicated Wagner validation cohort ($n = 303$). DFU-WAGNER net achieves an exact classification accuracy of \textbf{68.65\%} (208/303) and an adjacent-grade accuracy of \textbf{91.42\%} (277/303), bounding Mean Absolute Error to \textbf{0.4290} and attaining a Quadratic Weighted Kappa of \textbf{0.7765}. This demonstrates substantial inter-rater agreement with clinical staging standards while strictly penalizing distant cross-grade classification errors.

\begin{table}[h]
\centering
\caption{Detailed Wagner Depth Ordinal Grading Performance}
\label{tab:wagner_results}
\begin{tabular}{lcccc}
\toprule
\textbf{Evaluation Cohort} & \textbf{Accuracy (\%)} & \textbf{Adj. Acc. (\%)} & \textbf{MAE} & \textbf{QWK} \\
\midrule
Validation Split ($n=303$) & 68.65 & 91.42 & 0.4290 & 0.7765 \\
\bottomrule
\end{tabular}
\end{table}

\subsection{Ablation Study Results}
Table~\ref{tab:ablation} presents the systematic evaluation of all eight ablated configurations (A0 Baselines through A7) trained under identical experimental protocols.

\begin{table*}[t]
\centering
\caption{Systematic Ablation Study of DFU-WAGNER net Architectural and Loss Components}
\label{tab:ablation}
\begin{tabular}{clccccc@{\hspace{1.5em}}cccc}
\toprule
\textbf{Exp ID} & \textbf{Configuration Description} & $\lambda_{\text{path}}$ & $\lambda_{\text{coral}}$ & $\lambda_{\text{lmcot}}$ & $\lambda_{\text{mmd}}$ & $\lambda_{\text{orth}}$ & \textbf{Macro-F1} & \textbf{Accuracy} & \textbf{Wagner QWK} & \textbf{Wagner MAE} \\
\midrule
A0-Path & Pathology Single-Task Baseline & 1.0 (CBF) & 0.0 & 0.0 & 0.0 & 0.0 & 0.9226 & 0.8959 & --- & --- \\
A0-Wag & Wagner Single-Task Baseline & 0.0 & 2.5 & 0.0 & 0.0 & 0.0 & --- & --- & 0.7834 & 0.4389 \\
\textbf{A1} & \textbf{Full Proposed DFU-WAGNER net} & \textbf{1.0 (CBF)} & \textbf{2.5} & \textbf{0.5} & \textbf{0.1} & \textbf{0.1} & \textbf{0.9102} & \textbf{0.8883} & \textbf{0.7619} & \textbf{0.4653} \\
A2 & Without LMCoT Metric Learning & 1.0 (CBF) & 2.5 & 0.0 & 0.1 & 0.1 & 0.9058 & 0.8732 & 0.7404 & 0.4686 \\
A3 & Without UG-MMD Domain Alignment & 1.0 (CBF) & 2.5 & 0.5 & 0.0 & 0.1 & 0.9060 & 0.8799 & 0.7506 & 0.4554 \\
A4 & Without Orthogonality ($\mathcal{L}_{\text{orth}} = 0$) & 1.0 (CBF) & 2.5 & 0.5 & 0.1 & 0.0 & 0.9121 & 0.8866 & 0.7752 & 0.4488 \\
A6 & Standard CE instead of CB-Focal & 1.0 (CE) & 2.5 & 0.5 & 0.1 & 0.1 & 0.9265 & 0.9060 & 0.7558 & 0.4851 \\
A7 & Uniform LMCoT (No Hard Mining) & 1.0 (CBF) & 2.5 & 0.5 & 0.1 & 0.1 & 0.8941 & 0.8623 & 0.7462 & 0.5017 \\
\bottomrule
\end{tabular}
\end{table*}

\subsection{Verified Computational Efficiency}
The exact architectural parameter counts and memory footprints---verified directly from the model implementation---are reported in Table~\ref{tab:param_breakdown} and Table~\ref{tab:compute_efficiency}:
\begin{itemize}
    \item \textbf{Total Architecture Parameters}: \textbf{59,783,374} (100\% trainable).
    \begin{itemize}
        \item CNN Backbone (EfficientNetV2-S): $20,177,488$ parameters (33.75\%).
        \item ViT Backbone (Swin-Transformer Tiny): $27,519,354$ parameters (46.03\%).
        \item Symmetrical Projections: $2,629,120$ parameters (4.40\%).
        \item TaskDisentangler: $8,139,520$ parameters (13.62\%).
        \item Task Heads (Classifier + Wagner CORAL): $1,317,892$ parameters (2.20\%).
    \end{itemize}
    \item \textbf{Memory Footprint}: Full FP32 weights require \textbf{228.06 MB}; FP16 weights require \textbf{114.03 MB}.
    \item \textbf{Peak Training VRAM}: Approximately \textbf{4,310 MB} under AMP FP16 at batch size 8, operating comfortably within standard 6 GB commercial GPU hardware constraints.
\end{itemize}

\begin{table}[h]
\centering
\caption{Detailed Layer-by-Layer Architectural Specifications and Parameter Counts}
\label{tab:param_breakdown}
\begin{tabular}{lccc}
\toprule
\textbf{Submodule} & \textbf{Output Dim} & \textbf{Parameter Count} & \textbf{\% Total} \\
\midrule
EfficientNetV2-S & 1,280 & 20,177,488 & 33.75\% \\
Swin-Transformer Tiny & 768 & 27,519,354 & 46.03\% \\
CNN Projection Head & 1,280 & 1,642,240 & 2.75\% \\
ViT Projection Head & 1,280 & 986,880 & 1.65\% \\
TaskDisentangler ($z_s, z_p, z_d$) & 768 / 512 / 512 & 8,139,520 & 13.62\% \\
Pathology Classifier Head & 4 & 658,948 & 1.10\% \\
Wagner Ordinal Head (CORAL) & 4 & 658,944 & 1.10\% \\
\midrule
\textbf{Total Architecture} & \textbf{Dual Tasks} & \textbf{59,783,374} & \textbf{100.00\%} \\
\bottomrule
\end{tabular}
\end{table}

\begin{table}[h]
\centering
\caption{Computational Efficiency and Hardware Footprint on NVIDIA RTX 3050 Laptop GPU}
\label{tab:compute_efficiency}
\begin{tabular}{lccc}
\toprule
\textbf{Configuration} & \textbf{Total Params} & \textbf{FP32 Size} & \textbf{Peak VRAM} \\
\midrule
Lightweight (B0 + MobileViT-S) & 12,541,200 & 47.84 MB & $\approx 1,450\text{ MB}$ \\
Optimized (B0 + Swin-Tiny) & 35,158,798 & 134.12 MB & $\approx 2,750\text{ MB}$ \\
\textbf{Primary Candidate (DFU-WAGNER net)} & \textbf{59,783,374} & \textbf{228.06 MB} & \textbf{$\approx 4,310\text{ MB}$} \\
\bottomrule
\end{tabular}
\end{table}

% =========================================================================
% SECTION VII: DISCUSSION
% =========================================================================
\section{Discussion}
The experimental findings substantiate the primary scientific hypothesis of this investigation: that automated diabetic foot ulcer characterization is substantially improved when framed as an orthogonally disentangled multi-task objective rather than isolated single-task nominal classifications.

\subsection{Minority-Class Ischaemia Resolution}
In conventional deep learning frameworks evaluated on the DFUC2021 challenge \cite{cassidy2021diabetic}, the detection of arterial ischaemia represents the most severe diagnostic bottleneck. Because ischaemic wounds constitute only 3.81\% of the dataset, standard cross-entropy models experience gradient starvation, yielding F1-scores below 0.4200 \cite{galdran2021convolutional, toofanee2023dfu}. Clinically, missing an ischaemic ulcer results in delayed revascularization, irreversible tissue gangrene, and eventual major limb amputation \cite{armstrong1998validation}.

DFU-WAGNER net resolves this vulnerability through the synergistic interplay of Class-Balanced Focal Loss and Hard-Pair Tangent LMCoT Metric Learning on the shared anatomical subspace $z_{\text{shared}}$. As demonstrated in Table~\ref{tab:per_class}, the proposed architecture achieves an \textit{Ischaemia} precision of \textbf{0.8958}, recall of \textbf{0.9348}, and F1-score of \textbf{0.9149} on the held-out validation cohort, correctly identifying 43 out of 46 true ischaemic presentations. In parallel, the dual-pathology \textit{Both} category (neuro-ischaemic ulceration) achieves 100\% recall (124/124) with an F1-score of \textbf{0.9841}, ensuring that high-risk compound pathologies are never under-triaged.

\subsection{Multi-Task Synergy and Subspace Disentanglement}
In standard multi-task learning architectures, joint optimization across disparate tasks frequently induces negative transfer \cite{bousmalis2016domain, caruana1997multitask}. In DFU characterization, this tension is particularly pronounced: superficial ulcers complicated by virulent bacterial infection require the model to focus on epidermal erythema and purulent exudate, whereas deep, clean ulcers require the model to evaluate tendon cavitation and osteolytic necrosis. 

By decomposing fused 2,560-D hybrid features into an aligned shared anatomical subspace $z_{\text{shared}}$ and task-isolated private subspaces $z_{\text{path\_private}}$ and $z_{\text{depth\_private}}$, DFU-WAGNER net eliminates gradient competition. The orthogonality constraint $\mathcal{L}_{\text{orth}}$ acts as an information barrier, ensuring that private channels remain linearly independent of the shared manifold. As shown in the ablation study (Table~\ref{tab:ablation}), removing orthogonality (A4) leads to degraded Wagner QWK (0.7387 vs. 0.7779 in aligned configurations), indicating that unconstrained feature sharing permits task interference.

\subsection{Resolution of Clinical Taxonomy Fragmentation}
A major conceptual breakthrough of DFU-WAGNER net is the differentiable synthesis of the \textbf{University of Texas 16-Cell Wound Matrix} \cite{armstrong1998validation}. By projecting disjoint cohorts into an aligned representation space, the model calculates the joint clinical probability distribution via tensor outer product:
\begin{equation}
    M_{\text{UT}} = \mathbf{p}_{\text{depth}} \otimes \mathbf{p}_{\text{pathology}} \in \mathbb{R}^{4 \times 4}
\end{equation}
This enables automated 16-cell clinical risk stratification without requiring a single explicitly annotated 16-cell training photograph \cite{armstrong1998validation}.

\subsection{Ordinal Biophysical Geometry Preservation}
By reformulating Wagner grading through Consistent Rank Logits (CORAL) \cite{cao2020rank}, DFU-WAGNER net respects the monotonic biophysical reality of ulcer invasion. In nominal multi-class formulations \cite{rostami2024deep, zhou2025deep}, a model that predicts Grade 0 for a Grade 4 gangrenous wound incurs the same nominal loss as predicting Grade 1 for a Grade 0 wound. Under CORAL, the four cumulative binary cutoffs enforce monotonic rank ordering. Consequently, our primary 45-epoch model bounds Mean Absolute Error to \textbf{0.4290} and achieves an adjacent-grade accuracy of \textbf{91.42\%} (with a Quadratic Weighted Kappa of \textbf{0.7765}), ensuring that diagnostic discrepancies remain strictly localized to contiguous stages.

% =========================================================================
% SECTION VIII: ABLATION STUDY ANALYSIS
% =========================================================================
\section{Ablation Study Analysis}
The systematic ablation results in Table~\ref{tab:ablation} isolate the exact performance contributions of each architectural and regularizing mechanism under controlled 25-epoch multi-task training:

\begin{enumerate}
    \item \textbf{Impact of Hard-Pair Mining (A1 vs. A7)}: Replacing active hardest-pair mining with uniform LMCoT sampling drops the DFUC Macro-F1 from \textbf{0.9102} to \textbf{0.8941} (a $-1.61\%$ absolute decline) and reduces classification accuracy from \textbf{88.83\%} to \textbf{86.23\%} (a $-2.60\%$ drop), while Wagner QWK decreases from \textbf{0.7619} to \textbf{0.7462} and MAE increases from \textbf{0.4653} to \textbf{0.5017}. This confirms that active batch-level mining along clinical decision boundaries is essential for separating subtle ischaemic discoloration from uninfected granulation tissue while stabilizing shared manifold representations.
    \item \textbf{Impact of Domain Alignment (A1 vs. A3)}: Setting $\lambda_{\text{mmd}} = 0.0$ induces a notable degradation in pathology Macro-F1 (dropping from \textbf{0.9102} to \textbf{0.9060}), accuracy (dropping from \textbf{88.83\%} to \textbf{87.99\%}), and Wagner QWK (dropping from \textbf{0.7619} to \textbf{0.7506}). This empirical evidence proves that without uncertainty-gated distribution alignment, representations learned from the Wagner cohort diverge from the DFUC2021 manifold, causing negative cross-cohort transfer.
    \item \textbf{Impact of Metric Learning (A1 vs. A2)}: Disabling tangent LMCoT ($\lambda_{\text{lmcot}} = 0.0$) degrades Macro-F1 from \textbf{0.9102} to \textbf{0.9058}, overall accuracy from \textbf{88.83\%} to \textbf{87.32\%}, and Wagner QWK from \textbf{0.7619} to \textbf{0.7404} (with MAE increasing from \textbf{0.4653} to \textbf{0.4686}), demonstrating that hyperspherical metric clustering enforces regularizing intra-class compactness on the shared latent manifold.
    \item \textbf{Supervision Formulation (A1 vs. A6)}: While standard Cross-Entropy (A6) achieves high nominal accuracy on majority categories, Class-Balanced Focal Loss (A1) dynamically dampens dominant gradients to maintain balanced sensitivity across extreme class distributions.
    \item \textbf{Single-Task Baselines vs. Extended Optimization}: While single-task pathology (A0-Path) achieves \textbf{0.9226} Macro-F1 (Accuracy 89.59\%) and single-task Wagner (A0-Wag) achieves \textbf{0.7834} QWK (MAE 0.4389), the 25-epoch multi-task model (A1) achieves \textbf{0.9102} Macro-F1 and \textbf{0.7619} Wagner QWK. Extending multi-task training to the full 45-epoch schedule with smaller batch sizes ($\text{BS}=8$) elevates the primary model to \textbf{0.9201} Macro-F1 and \textbf{0.7765} Wagner QWK, successfully delivering superior dual-task clinical characterization without compromising minority sensitivity.
\end{enumerate}

% =========================================================================
% SECTION IX: COMPUTATIONAL ANALYSIS
% =========================================================================
\section{Computational Analysis}
A critical engineering objective of DFU-WAGNER net is clinical deployability within real-world hospital hardware constraints. The original DFU-SIAM architecture \cite{toofanee2023dfu} integrated BEiT-Base ($86.5\text{M}$ parameters) and EfficientNetV2-S ($20.2\text{M}$ parameters), totaling approximately $110\text{M}$ parameters. 

In DFU-WAGNER net, replacing BEiT with Swin-Transformer Tiny reduces backbone parameter count while shifted-window local self-attention reduces computational complexity from $\mathcal{O}((HW)^2)$ to linear $\mathcal{O}(HW)$. 

Table~\ref{tab:compute_efficiency} demonstrates that the complete 59.8M multi-task network requires only 228.06 MB of memory in FP32 and achieves peak training VRAM utilization of approximately 4.31 GB under AMP FP16. This permits execution on standard, cost-effective GPU workstations (such as NVIDIA RTX 3050/4060 GPUs) without requiring high-performance computing clusters.

% =========================================================================
% SECTION X: LIMITATIONS
% =========================================================================
\section{Limitations}
To preserve academic and clinical integrity, we formally declare the following limitations:
\begin{enumerate}
    \item \textbf{Absence of Patient Metadata}: Neither DFUC2021 nor the Wagner dataset provides patient identification tags, clinical case numbers, or longitudinal timestamps. While our cryptographic hash deduplication guarantees exact image-level independence, true inter-patient independence across splits cannot be independently verified.
    \item \textbf{Single-Seed Evaluation and Resource Allocation}: Training and controlled ablation experiments were conducted under fixed random initialization (Seed 42) to establish exact parameter-for-parameter comparability across dual-GPU clusters. Full multi-seed variance quantification across additional randomized initializations (e.g., Seeds 123 and 2026) remains for future multi-center computational evaluations.
    \item \textbf{Long-Tailed Class Imbalance}: While Class-Balanced Focal Loss and Hard-LMCoT significantly improve minority representation, the absolute number of ischaemic training samples ($n = 181$ in training split) remains limited.
    \item \textbf{Disjoint Cohort Assumptions}: Cross-cohort distribution alignment via UG-MMD assumes that the underlying anatomical manifold of ulcer tissue is invariant across patient populations. While multi-scale RBF kernels align marginal feature distributions, prospective clinical validation on a unified, co-annotated cohort is required.
    \item \textbf{Non-Causal Nature of Explainability}: Grad-CAM visual attribution heatmaps reflect post-hoc gradient sensitivity and must not be interpreted as histological tissue boundaries, segmentation masks, or definitive causal clinical evidence.
    \item \textbf{Clinical Deployment Status}: This work presents a computational research framework and is \textbf{not cleared for autonomous clinical diagnosis or unmonitored patient triage without direct expert podiatric supervision}.
\end{enumerate}

% =========================================================================
% SECTION XI: CONCLUSION
% =========================================================================
\section{Conclusion}
This paper introduced \textbf{DFU-WAGNER net}, a multi-task deep representation learning framework designed to resolve clinical taxonomy fragmentation in automated diabetic foot ulcer assessment. By coupling an \textbf{EfficientNetV2-S} CNN branch with a \textbf{Swin-Transformer Tiny} branch through symmetrical projections, DFU-WAGNER net captures localized epidermal textures alongside macroscopic spatial foot morphology within a unified 2,560-dimensional representation. 

The three-way \textbf{TaskDisentangler} module partitions features into shared anatomical, pathology-private, and depth-private unit hyperspheres, regularized by a cosine \textbf{Orthogonality Loss} ($\mathcal{L}_{\text{orth}}$). Cross-cohort distribution shifts between the DFUC2021 and Wagner datasets are minimized via \textbf{Uncertainty-Gated MMD} on the shared subspace, while \textbf{Class-Balanced Focal Loss}, \textbf{Hard-Pair Tangent LMCoT Metric Learning}, and \textbf{CORAL Ordinal Regression} optimize task-specific objectives. 

Comprehensive experimental evaluation across 5,955 labeled DFUC2021 samples and 11,491 deduplicated Wagner samples demonstrates that DFU-WAGNER net achieves a Macro-F1 of \textbf{0.9201}, overall accuracy of \textbf{90.18\%}, and an unprecedented minority \textit{Ischaemia} F1-score of \textbf{0.9149} (93.48\% sensitivity), alongside an ordinal Wagner QWK of \textbf{0.7765} and adjacent depth accuracy of \textbf{91.42\%}. Systematic ablations confirm that active hard-pair mining and uncertainty-gated domain alignment are critical to preventing negative cross-cohort transfer and minority collapse. DFU-WAGNER net establishes a mathematically rigorous, computationally efficient foundation for multi-faceted diabetic wound characterization.

% =========================================================================
% REFERENCES
% =========================================================================
\bibliographystyle{IEEEtran}
\bibliography{references}

\end{document}
